\clearpage                                  % В том числе гарантирует, что список литературы в оглавлении будет с правильным номером страницы
%\hypersetup{ urlcolor=black }               % Ссылки делаем чёрными
%\providecommand*{\BibDash}{}                % В стилях ugost2008 отключаем использование тире как разделителя
\urlstyle{rm}                               % ссылки URL обычным шрифтом
\ifdefmacro{\microtypesetup}{\microtypesetup{protrusion=false}}{} % не рекомендуется применять пакет микротипографики к автоматически генерируемому списку литературы
%\insertbibliofull                          % Все источники единым списком / all references in one list
% Режим с подсписками: сначала чужие работы, затем работы автора отдельным списком.
% Split mode: general literature first, then the author's own publications as a
% separate list (as in the dissertation).
\insertbiblioexternal                       % Литература (кроме работ автора) / general literature
% Для вывода работ автора выберите одно из двух / choose one for the author's works:
\insertbiblioauthor                         % работы автора единым списком / author's works in one list
%\insertbiblioauthorgrouped                 % работы автора по базам (ВАК, WoS, Scopus, ...) / grouped by index
\ifdefmacro{\microtypesetup}{\microtypesetup{protrusion=true}}{}
\urlstyle{tt}                               % возвращаем установки шрифта ссылок URL
%\hypersetup{ urlcolor={urlcolor} }          % Восстанавливаем цвет ссылок

% Подсчёт полного числа источников для счётчика \citenum (используется во введении)
% Count the total number of references for the \citenum counter (used in the intro).
\begingroup
\makeatletter
\defbibenvironment{justcount}
  {\setcounter{citenum}{0}}
  {\write\@auxout{%
    \string\gdef\string\totalsavedcitenum{\the\value{citenum}}%
  }}
  {\stepcounter{citenum}\renewcommand{\blx@driver}[1]{}}
\printbibliography[heading=nobibheading,env=justcount,keyword=bibliofull]
\makeatother
\endgroup
